\documentclass[twoside]{article}
\usepackage{mmpr}

\begin{document}

\Russian
\title{Классификация траекторий динамических систем с помощью физически-информированных нейронных сетей}
\author{Терентьев~А.\,А.}{Терентьев Александр Андреевич$^{1}$\speaker}{@}
\author{Стрижов~В.\,В.}{Стрижов Вадим Викторович$^{1}$}{@}
\organization{%
    $^1$Город, Организация}
\maketitle

В работе рассматривается задача классификации многомерных траекторий физических систем. Традиционные методы классификации временных рядов, в том числе глубинные модели, обычно опираются на статистические зависимости между соседними наблюдениями и плохо учитывают физическую природу процесса. Для реальных динамических систем такие подходы оказываются недостаточными, поскольку важную роль играют законы сохранения, уравнения движения и внутренние ограничения, определяющие поведение траектории.

Постановка задачи. Пусть $X=[\mathbf X_1,\dots,\mathbf X_T]$ --- траектория, где $\mathbf X_i=(\mathbf q_i,\dot{\mathbf q}_i)$, а $\mathbf y_i=\ddot{\mathbf q}_i$ --- ускорение. По обучающей выборке $D=\{(X_i,\mathbf y_i)\}_{i=1}^N$ требуется восстановить лагранжиан $L(\mathbf q,\dot{\mathbf q})$ и построить классификатор, который относит новую траекторию к одному из классов динамических систем.

Для восстановления модели используется физически-информированная лагранжева нейронная сеть. Ускорение определяется из уравнения Эйлера--Лагранжа:
$$
\ddot{\mathbf q}
=\left(\nabla_{\dot{\mathbf q}\dot{\mathbf q}}L\right)^{-1}
\left[\nabla_{\mathbf q}L-\left(\nabla_{\dot{\mathbf q}\mathbf q}L\right)\dot{\mathbf q}\right].
$$
Параметры сети обучаются по критерию
$$
\mathcal L(\mathbf w)=\frac{1}{T}\sum_{i=1}^{T}\|\hat{\mathbf y}_i-\mathbf y_i\|_2^2,
$$
где $\hat{\mathbf y}_i$ --- предсказанное ускорение. При этом решается задача оптимизации
$$
\mathbf w^* = \arg\min_{\mathbf w \in \mathbb W} \mathcal L(\mathbf w).
$$
В теоретической части работы вводится оператор
$$
A(L)=\frac{\partial L}{\partial \mathbf q}-\frac{\partial}{\partial \dot{\mathbf q}}\frac{\partial L}{\partial \dot{\mathbf q}}\dot{\mathbf q},
$$
при этом уравнение Эйлера--Лагранжа может быть переписано в операторной форме
$$
A(L)=\mathbf H(L)\ddot{\mathbf q},
\qquad
\mathbf H(L)=\nabla_{\dot{\mathbf q}\dot{\mathbf q}}L.
$$
Из-за инвариантности уравнения относительно умножения лагранжиана на положительную константу задача оптимизации допускает неоднозначность. Для регуляризации вводится ограничение
$$
\lambda(\mathbf H)>1,
$$
которое устраняет вырожденные решения. Однако этого недостаточно для полной однозначности, поэтому в качестве дополнительного слагаемого в функцию потерь используется мера несовместимости решения линейной системы
$$
\mathbf H\ddot{\mathbf q}=b,
\qquad
\mathbf H\hat{\ddot{\mathbf q}}=\hat b,
$$
где $\hat b$ определяется через производные лагранжиана. Тогда
$$
\mathcal L_{\text{add}}(\mathbf w)=\frac{1}{T}\sum_{i=1}^{T}\|\hat b_i-b_i\|_2^2.
$$
Согласно теоретическому результату, для точного восстановления ускорения выполняется $\mathcal L_{\text{add}}=0$. При этом, если выполнено ограничение $\lambda(\mathbf H)>1$, то минимизация полной функции потерь
$$
\mathcal L=\mathcal L_{\text{add}}+\mathcal L_{\text{restriction}}
$$
обеспечивает восстановление лагранжиана, для которого ускорение, вычисленное через уравнение Эйлера--Лагранжа, совпадает с наблюдаемым ускорением.
 
Основной результат работы состоит в том, что классификация выполняется не по самой траектории, а по восстановленному лагранжиану. Это позволяет учитывать физические ограничения системы и повышает устойчивость метода к нефизическим шумам. Предложенная метрика в пространстве лагранжианов достаточно чувствительна к различиям между классами движений и позволяет разделять траектории, порожденные разными физическими законами.

Важной частью метода является введение нормы в пространстве лагранжианов. Для каждой траектории строится соответствующий лагранжиан, а затем сравнение между различными системами выполняется не по самим траекториям, а по их лагранжианам. Такая метрика позволяет выделять классы движений по их физическому содержанию и делает классификацию устойчивой к случайным нефизическим возмущениям, которые могут возникать в данных. В работе также предполагается, что лагранжианы различных типов движений лежат в компактных и разделимых подмножествах, что делает возможным их разделение в пространстве моделей.

Экспериментальная часть выполнена на наборе PAMAP2, содержащем траектории человеческих движений, включая ходьбу, бег, приготовление пищи и другие повседневные действия. Полученные результаты показывают, что предложенный подход позволяет выделять физически различающиеся движения более надежно, чем методы, опирающиеся только на статистические свойства временных рядов. Это подтверждает целесообразность использования лагранжевой структуры и законов сохранения при классификации траекторий физических систем.

\begin{thebibliography}{1}
\bibitem{pinning}
    \emph{Raissi~M., Perdikaris~P., Karniadakis~G.~E.}
    Physics-informed neural networks: A deep learning framework for solving forward and inverse problems involving nonlinear partial differential equations~//
    Journal of Computational Physics.~--- 2019.~--- Vol. 378.~--- Pp. 686--707.
\bibitem{lnn}
    \emph{Cranmer~M., Greydanus~S., Haldane~D., et al.}
    Lagrangian neural networks~//
    Advances in Neural Information Processing Systems.~--- 2020.
\bibitem{pamap2}
    \emph{Reiss~A., Stricker~D.}
    Introducing a new benchmarked dataset for activity monitoring~//
    2012 16th International Symposium on Wearable Computers.~--- 2012.~--- Pp. 108--109.
\end{thebibliography}

\end{document}
